% =============================================================
%  ch14a_projekt0a_arm.tex  --  Projekt 0.A: Planarer 2-Gelenk-Arm
% =============================================================

\chapter{Projekt 0.A -- Planarer 2-Gelenk-Arm: Nodes, Topics, Services, Actions, IK}
\label{ch:projekt0a}

\textbf{Ziel:} ein erstes, vollständig selbst implementiertes ROS\,2-System, das
alle Grundbausteine berührt -- \textbf{Nodes}, \textbf{Topics},
\textbf{Services}, \textbf{Actions}, \textbf{URDF} und \textbf{inverse Kinematik}
-- und das wir in \textbf{RViz} visualisieren und in \textbf{Gazebo Harmonic}
simulieren. Wir übernehmen nur die \emph{Geometrie} eines bekannten
Lernroboters (den planaren 2-Gelenk-Arm im Stil des ``RRBot'' aus
\code{ros2\_control\_demos}) und schreiben die gesamte Logik selbst. Die
Sim-Real-Architektur aus Kapitel~\ref{ch:simreal} gilt unverändert.

\begin{infobox}[title=Warum gerade dieser Arm?]
Der Arm besteht aus zwei Drehgelenken (\emph{revolute}) und zwei Stäben und
bewegt sich in einer Ebene. Seine Geometrie ist trivial -- zwei Zylinder, ganz
ohne CAD-Mesh -- aber er deckt alle Lernziele ab. Wichtig: die \textbf{inverse
Kinematik eines 2-Glieder-Arms ist von Hand nachprüfbar} (es gibt sogar eine
geschlossene Lösung über den Kosinussatz). Wir nutzen den Arm aber, um das
\emph{allgemeine} IK-Verfahren zu üben, das auch bei komplexen Robotern greift:
Vorwärtskinematik aus \textbf{Denavit-Hartenberg}-Parametern plus
\textbf{numerische} Lösung über die Jacobi-Matrix. Die geschlossene Lösung dient
nur als Referenz zum Gegenprüfen. Vertiefung in Kapitel~\ref{ch:kinematik}.
\end{infobox}

\section{ROS\,2-Bausteine}
\begin{longtable}{@{}p{2.6cm}p{12cm}@{}}
\toprule
\textbf{Typ} & \textbf{Konkret (selbst implementiert, sofern nicht anders vermerkt)} \\
\midrule
Nodes    & \code{arm\_commander} (Action-/Service-Client) $\cdot$
           \code{reach\_action\_server} (IK + Bahnerzeugung) $\cdot$
           \code{robot\_state\_publisher} (Standard) $\cdot$
           \code{controller\_manager} + \code{joint\_trajectory\_controller} +
           \code{joint\_state\_broadcaster} (Standard) $\cdot$
           Hardware-Interface (Sim: \code{gz\_ros2\_control}) \\
Topics   & \code{/joint\_states} $\cdot$ \code{/tf}, \code{/tf\_static} $\cdot$
           \code{/robot\_description} $\cdot$
           \code{/arm\_controller/joint\_trajectory} \\
Services & \code{/arm/set\_named\_pose} (eigener \code{srv}: home/extended/folded) $\cdot$
           \code{/controller\_manager/switch\_controller} (Standard) \\
Actions  & \code{/arm/reach\_point} (eigene \code{action}: Ziel $x,y$;
           Feedback: Restdistanz; Result: Erfolg) $\cdot$ intern
           \code{control\_msgs/FollowJointTrajectory} \\
Parameter & Gliedlängen $l_1,l_2$ $\cdot$ Gelenklimits $\cdot$ benannte Posen \\
\bottomrule
\end{longtable}

\section{rqt\_graph (Zielbild)}
\begin{center}
\begin{tikzpicture}[node distance=6mm and 14mm]
\node[node] (cmd) {/arm\_commander};
\node[node,right=of cmd] (ras) {/reach\_action\_\\server\\\scriptsize (IK)};
\node[node,right=of ras] (jtc) {/joint\_trajectory\_\\controller};
\node[node,right=of jtc] (hw) {/gz\_ros2\_control\\\scriptsize (Sim-Hardware)};
\node[topic,below=10mm of jtc] (js) {/joint\_states};
\node[node,below=6mm of cmd] (rviz) {/rviz2};
\draw[flow] (cmd) -- node[topic,fill=white,inner sep=1pt]{\scriptsize reach\_point} (ras);
\draw[flow] (ras) -- node[topic,fill=white,inner sep=1pt]{\scriptsize traj goal} (jtc);
\draw[flow] (jtc) -- node[topic,fill=white,inner sep=1pt]{\scriptsize commands} (hw);
\draw[flow] (hw) |- (js);
\draw[flow] (js) -| (jtc);
\draw[flow] (js) -| (rviz);
\draw[flow,dashed,accent2] (ras) to[bend left=22]
  node[above,font=\scriptsize\itshape]{feedback} (cmd);
\end{tikzpicture}
\captionof{figure}{Projekt~0.A: Der Befehlsknoten schickt ein kartesisches Ziel
als Action-Goal; der Action-Server rechnet die IK, erzeugt eine Trajektorie und
übergibt sie dem Controller. \code{/joint\_states} schließt die Schleife und
speist RViz.}
\end{center}

\screenshotplaceholder{p0a-rqtgraph}{Tatsächlicher \code{rqt\_graph} des
laufenden Systems (\code{rqt\_graph} bzw.\ \code{ros2 run rqt\_graph
rqt\_graph}). Hier sieht man, ob alle Knoten und Topics wie geplant verbunden
sind -- das beste Werkzeug zur Fehlersuche in der Architektur.}

\section{Schritt für Schritt}

\subsection{Schritt 1 -- Pakete anlegen und die Struktur verstehen}
Die Projektlogik kommt in ein Python-Paket, die eigenen Schnittstellen
(\code{srv}/\code{action}) in ein separates \code{ament\_cmake}-Paket -- nur dort
generiert ROS\,2 aus \code{.srv}/\code{.action} echte Nachrichtentypen.
\begin{lstlisting}[style=bash,caption={Pakete anlegen (im Workspace-\code{src})}]
cd ~/ros2_ws/src
ros2 pkg create --build-type ament_python p0a_arm \
  --license Apache-2.0 \
  --dependencies rclpy sensor_msgs trajectory_msgs control_msgs

ros2 pkg create --build-type ament_cmake p0a_arm_interfaces \
  --license Apache-2.0 --dependencies rosidl_default_generators

mkdir -p p0a_arm/{urdf,config,launch}  p0a_arm_interfaces/{srv,action}
\end{lstlisting}

\paragraph{Was die Optionen bedeuten.}
\begin{itemize}
\item \code{--build-type ament\_python} -- das Paket wird über \code{setup.py}
  gebaut, die Knoten sind Python-Skripte. Alternative \code{ament\_cmake}
  (C++ oder Interface-Generierung).
\item \code{--dependencies \ldots} -- trägt die genannten Pakete als
  \code{<depend>} in die \code{package.xml} ein. \code{rclpy} ist die
  Python-Client-Bibliothek von ROS\,2 (Knoten, Publisher, Timer \ldots), dazu
  die verwendeten Nachrichtenpakete. \emph{Deklaration}, keine Installation.
\item \code{--license Apache-2.0} -- vermeidet die TODO-Warnung und schreibt
  eine gültige Lizenz in \code{package.xml}.
\end{itemize}

\paragraph{Zielstruktur des Python-Pakets.}
\begin{lstlisting}[style=bash,caption={Verzeichnisbaum von \code{p0a\_arm} (Endstand Projekt 0.A)}]
p0a_arm/
|- package.xml          # Metadaten + Abhaengigkeiten (<depend>)
|- setup.py             # Build/Install: Knoten (entry_points) + Daten (data_files)
|- setup.cfg            # Ablageort der ausfuehrbaren Skripte
|- resource/p0a_arm     # Marker fuer den ament-Index
|- p0a_arm/             # Python-Modul (gleicher Name) -> hier die Knoten
|  |- __init__.py
|  |- arm_commander.py
|  |- reach_action_server.py
|  `- ik.py
|- urdf/    arm.urdf.xacro
|- config/  controllers.yaml  arm.rviz
`- launch/  display.launch.py  gazebo.launch.py
\end{lstlisting}
\noindent Die \code{package.xml} ist die Visitenkarte des Pakets: Name, Version,
Lizenz und die \code{<depend>}-Zeilen, die \code{rosdep} und \code{colcon}
auswerten. Die \code{setup.py} bestimmt, \emph{was installiert wird} -- dazu in
Schritt~3 mehr: \code{urdf/}, \code{launch/} und \code{config/} müssen dort
explizit eingetragen werden, sonst landen sie nicht im \code{install/}-Baum.
\begin{lstlisting}[style=bash,caption={\code{ros2\_ws}-Delta nach Schritt 1 (\# neu / \# angepasst)}]
src/
|- p0a_arm/             # neu (ament_python: rclpy-Knoten)
`- p0a_arm_interfaces/  # neu (ament_cmake: srv/action)
\end{lstlisting}

\subsection{Schritt 2 -- Geometrie als URDF/Xacro}
Wir schreiben das Modell selbst -- das \emph{ist} das URDF-Lernziel. Datei:
\code{p0a\_arm/urdf/arm.urdf.xacro}. Beide Gelenke drehen um $z$, die Glieder
liegen entlang $+x$ ($l_1=1{,}0$, $l_2=0{,}8$, passend zur DH-Tabelle aus
Schritt~7).

\paragraph{Aufbau einer URDF.} Eine URDF beschreibt einen Baum aus
\textbf{Links} (starre Körper) und \textbf{Joints} (Verbindungen):
\begin{itemize}
\item \textbf{\code{<robot name=\ldots>}} -- Wurzelelement, allgemeine Metadaten.
  Der Robotername ist frei; entscheidend sind die später referenzierten
  \emph{Link-} und \emph{Joint-Namen}.
\item \textbf{\code{<joint>}} -- verbindet einen \code{<parent>}- mit einem
  \code{<child>}-Link, mit \code{<origin>} (Lage relativ zum Parent),
  \code{<axis>} (Achse) und \code{<limit>} (Grenzen). \code{type}:
  \code{revolute} (begrenzt drehbar), \code{continuous}, \code{prismatic}
  (schiebend), \code{fixed} (starr).
\item \textbf{\code{<link>}} -- ein Körper mit \code{<visual>} (Darstellung),
  \code{<collision>} (Kollisionskörper) und \code{<inertial>} (Masse +
  Trägheitstensor; für Gazebo nötig, für RViz optional).
\item \textbf{\code{<ros2\_control>}/\code{<gazebo>}} -- keine Geometrie, sondern
  die Anbindung an die Simulation (Schritt~4).
\end{itemize}

\begin{infobox}[title=Müssen Joints und Links in Reihenfolge stehen?]
\textbf{Nein.} URDF wird \emph{nicht} als Reihenfolge von oben nach unten
gelesen. Der kinematische Baum entsteht allein aus den
\code{<parent>}/\code{<child>}-Verweisen der Joints. Die Abfolge \code{joint1},
\code{link1}, \code{joint2}, \code{link2} \ldots\ ist nur \emph{lesbare}
Konvention -- man dürfte alle Links zuerst und dann alle Joints notieren. Pflicht
ist nur: genau \emph{ein} Wurzel-Link (ohne Parent), ein zusammenhängender,
zyklenfreier Baum, und jeder referenzierte Link existiert.
\end{infobox}

\noindent Die vollständige, lauffähige Datei (Glieder als Boxen, mit
Inertials und Collisions):
\begin{lstlisting}[style=xml,caption={\code{urdf/arm.urdf.xacro} -- vollständig}]
<?xml version="1.0"?>
<!-- Planarer 2-Gelenk-Arm: Gelenke um z, Glieder entlang +x. l1=1.0, l2=0.8 -->
<robot name="rrarm">

  <link name="world"/>
  <joint name="fix" type="fixed">
    <parent link="world"/><child link="base"/>
  </joint>

  <link name="base">
    <visual><origin xyz="0 0 0.05"/>
      <geometry><cylinder length="0.1" radius="0.05"/></geometry>
      <material name="grey"><color rgba="0.5 0.5 0.5 1"/></material></visual>
    <collision><origin xyz="0 0 0.05"/>
      <geometry><cylinder length="0.1" radius="0.05"/></geometry></collision>
    <inertial><mass value="1.0"/>
      <inertia ixx="0.01" ixy="0" ixz="0" iyy="0.01" iyz="0" izz="0.01"/></inertial>
  </link>

  <joint name="joint1" type="revolute">
    <parent link="base"/><child link="link1"/>
    <origin xyz="0 0 0.1"/><axis xyz="0 0 1"/>
    <limit lower="-3.14" upper="3.14" effort="50" velocity="2.0"/>
  </joint>
  <link name="link1">
    <visual><origin xyz="0.5 0 0"/>
      <geometry><box size="1.0 0.06 0.06"/></geometry>
      <material name="blue"><color rgba="0.2 0.4 0.8 1"/></material></visual>
    <collision><origin xyz="0.5 0 0"/>
      <geometry><box size="1.0 0.06 0.06"/></geometry></collision>
    <inertial><origin xyz="0.5 0 0"/><mass value="1.0"/>
      <inertia ixx="0.001" ixy="0" ixz="0" iyy="0.084" iyz="0" izz="0.084"/></inertial>
  </link>

  <joint name="joint2" type="revolute">
    <parent link="link1"/><child link="link2"/>
    <origin xyz="1.0 0 0"/><axis xyz="0 0 1"/>
    <limit lower="-3.14" upper="3.14" effort="50" velocity="2.0"/>
  </joint>
  <link name="link2">
    <visual><origin xyz="0.4 0 0"/>
      <geometry><box size="0.8 0.05 0.05"/></geometry>
      <material name="orange"><color rgba="0.9 0.5 0.1 1"/></material></visual>
    <collision><origin xyz="0.4 0 0"/>
      <geometry><box size="0.8 0.05 0.05"/></geometry></collision>
    <inertial><origin xyz="0.4 0 0"/><mass value="0.8"/>
      <inertia ixx="0.0008" ixy="0" ixz="0" iyy="0.043" iyz="0" izz="0.043"/></inertial>
  </link>

  <!-- Anbindung an die Simulation (Schritt 4) -->
  <ros2_control name="GazeboSystem" type="system">
    <hardware><plugin>gz_ros2_control/GazeboSimSystem</plugin></hardware>
    <joint name="joint1"><command_interface name="position"/>
      <state_interface name="position"/><state_interface name="velocity"/></joint>
    <joint name="joint2"><command_interface name="position"/>
      <state_interface name="position"/><state_interface name="velocity"/></joint>
  </ros2_control>
</robot>
\end{lstlisting}

\begin{warnbox}[title=Der \code{<gazebo>}-Plugin-Block kommt erst in Schritt 4]
Der Block, der die Controller lädt, verweist mit \code{\$(find p0a\_arm)} auf das
\emph{installierte} Paket:
\begin{lstlisting}[style=xml]
<gazebo>
  <plugin filename="gz_ros2_control-system"
          name="gz_ros2_control::GazeboSimROS2ControlPlugin">
    <parameters>$(find p0a_arm)/config/controllers.yaml</parameters>
  </plugin>
</gazebo>
\end{lstlisting}
\code{\$(find \ldots)} lässt sich erst auflösen, \emph{nachdem} das Paket gebaut
und gesourct ist. Schon für den RViz-Test (Schritt~3) hineingenommen, scheitert
\code{xacro} mit ``package not found''. Daher erst in Schritt~4 ergänzen.
\end{warnbox}
\begin{lstlisting}[style=bash,caption={\code{ros2\_ws}-Delta nach Schritt 2 (\# neu / \# angepasst)}]
src/p0a_arm/
`- urdf/arm.urdf.xacro  # neu (Robotermodell: links, joints, ros2_control)
\end{lstlisting}

\subsection{Schritt 3 -- In RViz darstellen (über eine Launch-Datei)}
Den großen URDF-String \emph{nicht} als Kommandozeilen-Parameter übergeben --
\code{rcl} kann das mehrzeilige XML nicht parsen (\code{Couldn't parse parameter
override rule}). Der saubere Weg ist eine \textbf{Launch-Datei}, die
\code{robot\_state\_publisher} den \code{robot\_description}-Parameter als echten
Wert übergibt.

\paragraph{Ressourcen installierbar machen.} Damit \code{ros2 launch} die Dateien
findet, müssen \code{urdf/}, \code{launch/} und \code{config/} in der
\code{setup.py} als \code{data\_files} eingetragen werden:
\begin{lstlisting}[style=py,caption={Ergänzung in \code{setup.py}}]
import os
from glob import glob
# ... in data_files=[ ... ] zusaetzlich:
(os.path.join('share', package_name, 'urdf'),   glob('urdf/*')),
(os.path.join('share', package_name, 'launch'), glob('launch/*.launch.py')),
(os.path.join('share', package_name, 'config'), glob('config/*')),
\end{lstlisting}

\begin{lstlisting}[style=py,caption={\code{launch/display.launch.py}}]
from launch import LaunchDescription
from launch.substitutions import Command, PathJoinSubstitution
from launch_ros.actions import Node
from launch_ros.substitutions import FindPackageShare
from launch_ros.parameter_descriptions import ParameterValue

def generate_launch_description():
    pkg = FindPackageShare('p0a_arm')
    xacro_file = PathJoinSubstitution([pkg, 'urdf', 'arm.urdf.xacro'])
    rviz_config = PathJoinSubstitution([pkg, 'config', 'arm.rviz'])
    robot_description = ParameterValue(Command(['xacro ', xacro_file]),
                                       value_type=str)
    return LaunchDescription([
        Node(package='robot_state_publisher', executable='robot_state_publisher',
             parameters=[{'robot_description': robot_description}]),
        Node(package='joint_state_publisher_gui',
             executable='joint_state_publisher_gui'),
        Node(package='rviz2', executable='rviz2',
             arguments=['-d', rviz_config]),
    ])
\end{lstlisting}

\paragraph{RViz-Ansicht vorkonfigurieren (automatisch).} Frisch gestartet steht
RViz auf \emph{Fixed Frame} \code{map} (Fehler ``Frame [map] does not exist'')
und zeigt kein Modell -- es ist noch kein RobotModel-Display angelegt. Damit
Modell und TF-Achsen sofort erscheinen, legt die mitgelieferte
\code{config/arm.rviz} drei Dinge fest:
\begin{lstlisting}[style=yaml,caption={\code{config/arm.rviz} -- die entscheidenden Einträge}]
Visualization Manager:
  Global Options:
    Fixed Frame: world             # statt "map" -> kein Frame-Fehler mehr
  Displays:
    - Class: rviz_default_plugins/RobotModel
      Description Source: Topic
      Description Topic:
        Value: /robot_description
        Durability Policy: Transient Local   # WICHTIG: rsp latcht die Beschreibung
    - Class: rviz_default_plugins/TF           # zeigt die Gelenk-Frames
\end{lstlisting}

\paragraph{Der manuelle Weg (ohne Config-Datei).} In \emph{Global Options} das
\emph{Fixed Frame} auf \code{world} setzen. Dann unten links \emph{Add}
$\rightarrow$ \code{RobotModel} $\rightarrow$ \emph{OK} und in dessen
Eigenschaften \emph{Description Topic} auf \code{/robot\_description} stellen --
der Arm erscheint. Optional \emph{Add} $\rightarrow$ \code{TF} für die
Gelenk-Frames. Zum Schluss \emph{File\,$\rightarrow$\,Save Config As} nach
\code{config/arm.rviz}, damit der \code{-d}-Start sie künftig automatisch lädt.
\textbf{Häufige Falle:} ohne \emph{Transient Local} sieht das RobotModel die
einmalig gelatchte Beschreibung nicht -- das Modell bleibt unsichtbar.

\begin{lstlisting}[style=bash,caption={Bauen, sourcen, starten}]
cd ~/ros2_ws
colcon build --packages-select p0a_arm
source install/setup.bash          # neue Terminals tun das via .bashrc selbst
ros2 launch p0a_arm display.launch.py   # laedt arm.rviz automatisch (-d)
\end{lstlisting}

\begin{lstlisting}[style=bash,caption={\code{ros2\_ws}-Delta nach Schritt 3 (\# neu / \# angepasst)}]
src/p0a_arm/
|- setup.py                   # angepasst: data_files (urdf, launch, config)
|- launch/display.launch.py   # neu
`- config/arm.rviz            # neu
\end{lstlisting}

\begin{figure}[htbp]\centering
  \includegraphics[width=0.85\textwidth]{{project_0.A/step3_rviz_with_state_publisher}.png}
  \caption{RViz mit dem Arm-Modell und den TF-Achsen. Über den
  \code{joint\_state\_publisher\_gui} lassen sich die Gelenke bewegen -- dreht
  sich der Arm plausibel, sind URDF und TF korrekt.}
  \label{fig:p0a-rviz-model}
\end{figure}

\subsection{Schritt 4 -- In Gazebo simulieren (\code{ros2\_control})}
Jetzt mit Physik. \textbf{Zuerst} den in Schritt~2 gezeigten
\code{<gazebo>}-Block in \code{arm.urdf.xacro} aktivieren (Auskommentierung
aufheben): Da das Paket nun gebaut und gesourct ist, löst \code{\$(find
p0a\_arm)} auf den installierten \code{config}-Pfad auf. Die Controller
beschreibt eine YAML, die der vom Plugin gestartete \code{controller\_manager}
lädt:
\begin{lstlisting}[style=yaml,caption={\code{config/controllers.yaml}}]
controller_manager:
  ros__parameters:
    update_rate: 100
    joint_state_broadcaster:
      type: joint_state_broadcaster/JointStateBroadcaster
    arm_controller:
      type: joint_trajectory_controller/JointTrajectoryController
arm_controller:
  ros__parameters:
    joints: [joint1, joint2]
    command_interfaces: [position]
    state_interfaces: [position, velocity]
\end{lstlisting}
\noindent Nach jeder Änderung an \code{urdf/} oder \code{config/} neu bauen und
sourcen, sonst läuft die alte installierte Fassung:
\code{colcon build --packages-select p0a\_arm \&\& source install/setup.bash}.

\paragraph{Empfohlen: alles in einer Launch-Datei.} Vier Knoten von Hand in
mehreren Terminals zu starten ist fehleranfällig (Gazebo blockiert sein Terminal,
\code{robot\_state\_publisher} muss laufen, die Controller dürfen erst nach dem
Spawn kommen). Eine Launch-Datei erledigt das in \emph{einem} Befehl und
sequenziert die Schritte über Event-Handler:
\begin{lstlisting}[style=py,caption={\code{launch/gazebo.launch.py}}]
from launch import LaunchDescription
from launch.actions import IncludeLaunchDescription, RegisterEventHandler
from launch.event_handlers import OnProcessExit
from launch.launch_description_sources import PythonLaunchDescriptionSource
from launch.substitutions import Command, PathJoinSubstitution
from launch_ros.actions import Node
from launch_ros.substitutions import FindPackageShare
from launch_ros.parameter_descriptions import ParameterValue

def generate_launch_description():
    pkg = FindPackageShare('p0a_arm')
    ros_gz_sim = FindPackageShare('ros_gz_sim')
    xacro_file = PathJoinSubstitution([pkg, 'urdf', 'arm.urdf.xacro'])
    robot_description = ParameterValue(Command(['xacro ', xacro_file]),
                                       value_type=str)
    rsp = Node(package='robot_state_publisher', executable='robot_state_publisher',
               parameters=[{'robot_description': robot_description,
                            'use_sim_time': True}])   # TF-Stamps = Gazebo-Zeit
    gz = IncludeLaunchDescription(
        PythonLaunchDescriptionSource(
            PathJoinSubstitution([ros_gz_sim, 'launch', 'gz_sim.launch.py'])),
        launch_arguments={'gz_args': '-r empty.sdf'}.items())   # -r: sofort laufen
    # /clock von Gazebo nach ROS bruecken ('[' = gz -> ROS) -- Zeitquelle fuer
    # alle use_sim_time-Knoten (controller_manager, rsp, RViz)
    clock_bridge = Node(package='ros_gz_bridge', executable='parameter_bridge',
                        arguments=['/clock@rosgraph_msgs/msg/Clock[gz.msgs.Clock'])
    spawn = Node(package='ros_gz_sim', executable='create',
                 arguments=['-topic', 'robot_description', '-name', 'rrarm'])
    jsb = Node(package='controller_manager', executable='spawner',
               arguments=['joint_state_broadcaster'])
    arm = Node(package='controller_manager', executable='spawner',
               arguments=['arm_controller'])
    return LaunchDescription([
        rsp, gz, clock_bridge, spawn,
        RegisterEventHandler(OnProcessExit(target_action=spawn, on_exit=[jsb])),
        RegisterEventHandler(OnProcessExit(target_action=jsb, on_exit=[arm])),
    ])
\end{lstlisting}

\noindent Damit genügen \textbf{zwei Terminals}: eines startet alles, das zweite
prüft die Controller und schickt eine Trajektorie.
\begin{lstlisting}[style=bash,caption={Terminal 1 -- bauen, sourcen, alles starten}]
cd ~/ros2_ws
colcon build --packages-select p0a_arm
source install/setup.bash
ros2 launch p0a_arm gazebo.launch.py
\end{lstlisting}
\begin{lstlisting}[style=bash,caption={Terminal 2 -- Controller prüfen und Arm kommandieren}]
source ~/ros2_ws/install/setup.bash
ros2 control list_controllers          # beide sollten "active" sein
ros2 topic pub --once /arm_controller/joint_trajectory \
  trajectory_msgs/msg/JointTrajectory \
  '{joint_names: [joint1, joint2],
    points: [{positions: [0.8, -0.6], time_from_start: {sec: 2}}]}'
\end{lstlisting}

\paragraph{Was dabei unter der Haube passiert.} Die Launch-Datei ersetzt diese
vier Einzelschritte -- nützlich zum Verständnis, von Hand aber nicht nötig:
\begin{lstlisting}[style=bash,caption={Äquivalente Einzelbefehle}]
ros2 launch ros_gz_sim gz_sim.launch.py gz_args:="-r empty.sdf"
ros2 run ros_gz_sim create -topic robot_description -name rrarm
ros2 control load_controller --set-state active joint_state_broadcaster
ros2 control load_controller --set-state active arm_controller
\end{lstlisting}
\paragraph{Was die Befehle und Flags tun.}
\begin{itemize}
\item \code{ros\_gz\_sim gz\_sim.launch.py gz\_args:="empty.sdf"} -- startet
  Gazebo Harmonic mit einer leeren Welt.
\item \code{ros\_gz\_sim create -topic robot\_description} -- spawnt das Modell
  aus dem (vom \code{robot\_state\_publisher}) veröffentlichten
  \code{robot\_description}-Topic in die laufende Welt.
\item \code{ros2 control load\_controller --set-state active <name>} -- lädt
  einen Controller \emph{und} schaltet ihn in einem Schritt aktiv. Ohne
  \code{--set-state active} bliebe er nur konfiguriert (inaktiv) und würde keine
  Befehle umsetzen.
\item \code{controller\_manager spawner <name>} -- die Launch-Variante davon:
  lädt, konfiguriert und aktiviert einen Controller in einem Schritt.
\item \textbf{Reihenfolge:} erst \code{joint\_state\_broadcaster} (publiziert
  \code{/joint\_states}), dann \code{arm\_controller} (nimmt Trajektorien an).
\end{itemize}
\begin{lstlisting}[style=bash,caption={\code{ros2\_ws}-Delta nach Schritt 4 (\# neu / \# angepasst)}]
src/p0a_arm/
|- urdf/arm.urdf.xacro       # angepasst: <gazebo>-Plugin-Block aktiviert
|- config/controllers.yaml   # neu
`- launch/gazebo.launch.py   # neu (Gazebo + rsp + spawn + Controller)
\end{lstlisting}

\begin{figure}[htbp]\centering
  \includegraphics[width=0.85\textwidth]{{project_0.A/step4_gazebo}.png}
  \caption{Der Arm in Gazebo Harmonic. Hier wirkt Schwerkraft: ohne aktiven
  Controller sinkt der Arm -- mit aktivem \code{arm\_controller} hält er die per
  Trajektorie kommandierte Stellung.}
  \label{fig:p0a-gazebo}
\end{figure}

\subsection{Schritt 5 -- Topics: kommandieren und auslesen}
Erst direkt von der Kommandozeile, dann aus einem eigenen Knoten:
\begin{lstlisting}[style=bash,caption={Gelenke per Topic ansprechen}]
ros2 topic echo /joint_states                       # Ist-Zustand mitlesen
ros2 topic pub --once /arm_controller/joint_trajectory \
  trajectory_msgs/msg/JointTrajectory \
  '{joint_names: [joint1, joint2],
    points: [{positions: [0.5, -0.8], time_from_start: {sec: 2}}]}'
\end{lstlisting}
\noindent Als Knoten verpackt fährt \code{arm\_commander} abwechselnd zwei Posen
an, sodass eine sichtbare Bewegung entsteht.
\begin{lstlisting}[style=py,caption={\code{p0a\_arm/arm\_commander.py}}]
import rclpy
from rclpy.node import Node
from builtin_interfaces.msg import Duration
from trajectory_msgs.msg import JointTrajectory, JointTrajectoryPoint

class ArmCommander(Node):
    def __init__(self):
        super().__init__('arm_commander')
        self.pub = self.create_publisher(
            JointTrajectory, '/arm_controller/joint_trajectory', 10)
        self.poses = [[0.8, -0.6], [-0.8, 0.6]]    # zwei Zielposen
        self.idx = 0
        self.create_timer(3.0, self.tick)           # alle 3 s neues Ziel

    def tick(self):
        target = self.poses[self.idx]
        self.idx = (self.idx + 1) % len(self.poses)
        msg = JointTrajectory()
        msg.joint_names = ['joint1', 'joint2']
        p = JointTrajectoryPoint()
        p.positions = target
        p.time_from_start = Duration(sec=2)         # 2 s bis zum Ziel
        msg.points = [p]
        self.pub.publish(msg)

def main():
    rclpy.init()
    rclpy.spin(ArmCommander())
    rclpy.shutdown()
\end{lstlisting}
\noindent Den Knoten in \code{setup.py} als Skript registrieren (sonst kennt
\code{ros2 run} ihn nicht):
\begin{lstlisting}[style=py,caption={\code{p0a\_arm/setup.py} -- \code{entry\_points}}]
entry_points={
    'console_scripts': [
        'arm_commander = p0a_arm.arm_commander:main',
    ],
},
\end{lstlisting}

\paragraph{Ausführen (drei Terminals).} Damit sich etwas bewegt, muss der
Gazebo-Stack aus Schritt~4 laufen -- er stellt den \code{arm\_controller}, der
die Trajektorien ausführt.
\begin{lstlisting}[style=bash,caption={Terminal 1 -- Gazebo-Stack (wie Schritt 4)}]
cd ~/ros2_ws
colcon build --packages-select p0a_arm
source install/setup.bash
ros2 launch p0a_arm gazebo.launch.py
\end{lstlisting}
\begin{lstlisting}[style=bash,caption={Terminal 2 -- den Knoten starten}]
source ~/ros2_ws/install/setup.bash
ros2 run p0a_arm arm_commander
\end{lstlisting}
\begin{lstlisting}[style=bash,caption={Terminal 3 (optional) -- RViz eigenständig}]
source ~/ros2_ws/install/setup.bash
rviz2 -d "$(ros2 pkg prefix --share p0a_arm)/config/arm.rviz" \
  --ros-args -p use_sim_time:=true     # mit Gazebo: sonst "No transform"
\end{lstlisting}

\begin{warnbox}[title={RViz neben Gazebo -- nicht über display.launch.py}]
Der Arm bewegt sich bereits sichtbar in \textbf{Gazebo}. Will man ihn
\emph{zusätzlich} in RViz sehen, RViz \textbf{eigenständig} starten (Terminal 3)
-- \emph{nicht} über \code{display.launch.py}: dessen
\code{joint\_state\_publisher\_gui} würde mit dem Gazebo-%
\code{joint\_state\_broadcaster} um \code{/joint\_states} konkurrieren. RViz
zeigt die Bewegung dann über \code{/joint\_states}\,$\rightarrow$\,TF.
\end{warnbox}
\begin{lstlisting}[style=bash,caption={\code{ros2\_ws}-Delta nach Schritt 5 (\# neu / \# angepasst)}]
src/p0a_arm/
|- setup.py                  # angepasst: entry_points -> arm_commander
`- p0a_arm/arm_commander.py  # neu
\end{lstlisting}

\subsection{Schritt 6 -- Service: in eine benannte Pose fahren}
Ein Service ist die richtige Wahl für einen kurzen Befehl mit klarer Antwort.
\begin{lstlisting}[style=bash,caption={\code{srv/SetNamedPose.srv}}]
string name        # "home" | "extended" | "folded"
---
bool ok
string message
\end{lstlisting}
\paragraph{Interface bauen.} Eigene \code{srv}-Typen entstehen nur in einem
\code{ament\_cmake}-Paket. In \code{p0a\_arm\_interfaces} zwei Dinge ergänzen --
die Generierung in \code{CMakeLists.txt} und die Laufzeit-Abhängigkeiten in
\code{package.xml}:
\begin{lstlisting}[style=bash,caption={\code{p0a\_arm\_interfaces/CMakeLists.txt} (Ergänzung)}]
find_package(rosidl_default_generators REQUIRED)

rosidl_generate_interfaces(${PROJECT_NAME}
  "srv/SetNamedPose.srv"
)
ament_export_dependencies(rosidl_default_runtime)
\end{lstlisting}
\begin{lstlisting}[style=xml,caption={\code{p0a\_arm\_interfaces/package.xml} (Ergänzung)}]
<exec_depend>rosidl_default_runtime</exec_depend>
<!-- am Ende, direkt vor <export>: -->
<member_of_group>rosidl_interface_packages</member_of_group>
\end{lstlisting}

\paragraph{Server-Knoten.} Er bildet den Namen auf Gelenkwinkel ab, sendet die
Trajektorie (wie in Schritt~5) und meldet Erfolg in der Antwort.
\begin{lstlisting}[style=py,caption={\code{p0a\_arm/named\_pose\_server.py}}]
import rclpy
from rclpy.node import Node
from builtin_interfaces.msg import Duration
from trajectory_msgs.msg import JointTrajectory, JointTrajectoryPoint
from p0a_arm_interfaces.srv import SetNamedPose

class NamedPoseServer(Node):
    NAMED_POSES = {
        'extended': [0.0, 0.0],    # voll gestreckt
        'home':     [1.57, 0.0],   # 90-Grad-Ruhestellung
        'folded':   [0.0, 2.8],    # link2 zurueckgeklappt
    }
    def __init__(self):
        super().__init__('named_pose_server')
        self.pub = self.create_publisher(
            JointTrajectory, '/arm_controller/joint_trajectory', 10)
        self.create_service(SetNamedPose, '/arm/set_named_pose', self.on_request)

    def on_request(self, request, response):
        target = self.NAMED_POSES.get(request.name)
        if target is None:
            response.ok = False
            response.message = f"unbekannte Pose '{request.name}'"
            return response
        msg = JointTrajectory()
        msg.joint_names = ['joint1', 'joint2']
        p = JointTrajectoryPoint()
        p.positions = target
        p.time_from_start = Duration(sec=2)
        msg.points = [p]
        self.pub.publish(msg)
        response.ok = True
        response.message = f"fahre zu '{request.name}' -> {target}"
        return response

def main():
    rclpy.init()
    rclpy.spin(NamedPoseServer())
    rclpy.shutdown()
\end{lstlisting}
\begin{warnbox}[title={Callback nicht \code{handle} nennen}]
Die Callback-Methode heißt bewusst \code{on\_request} und \emph{nicht}
\code{handle}: \code{rclpy.node.Node} besitzt bereits eine Eigenschaft
\code{handle} (den internen rcl-Knoten-Handle). Eine eigene Methode \code{handle}
überschreibt sie, und schon \code{super().\_\_init\_\_()} scheitert mit
\code{'method' object does not support the context manager protocol}. Gleiches
gilt für andere Node-Member (\code{context}, \code{executor}, \ldots).
\end{warnbox}
\noindent Im \code{p0a\_arm}-Paket die Abhängigkeit (für den \code{import}) und
den Einstiegspunkt eintragen:
\begin{lstlisting}[style=xml,caption={\code{p0a\_arm/package.xml}}]
<depend>p0a_arm_interfaces</depend>
\end{lstlisting}
\begin{lstlisting}[style=py,caption={\code{p0a\_arm/setup.py} -- \code{entry\_points}}]
'named_pose_server = p0a_arm.named_pose_server:main',
\end{lstlisting}

\paragraph{Bauen und ausführen.} Das Interface-Paket muss \emph{vor}
\code{p0a\_arm} gebaut werden; \code{colcon} ordnet das anhand der
\code{<depend>}-Zeile automatisch.
\begin{lstlisting}[style=bash,caption={Terminal 1 -- bauen + Gazebo-Stack}]
cd ~/ros2_ws
colcon build --packages-select p0a_arm_interfaces p0a_arm
source install/setup.bash
ros2 launch p0a_arm gazebo.launch.py
\end{lstlisting}
\begin{lstlisting}[style=bash,caption={Terminal 2 -- Server bereitstellen}]
source ~/ros2_ws/install/setup.bash
ros2 run p0a_arm named_pose_server
\end{lstlisting}
\begin{lstlisting}[style=bash,caption={Terminal 3 -- Service aufrufen}]
source ~/ros2_ws/install/setup.bash
ros2 service call /arm/set_named_pose \
  p0a_arm_interfaces/srv/SetNamedPose "{name: folded}"
ros2 service call /arm/set_named_pose \
  p0a_arm_interfaces/srv/SetNamedPose "{name: extended}"
ros2 service call /arm/set_named_pose \
  p0a_arm_interfaces/srv/SetNamedPose "{name: home}"
\end{lstlisting}

\begin{lstlisting}[style=bash,caption={\code{ros2\_ws}-Delta nach Schritt 6 (\# neu / \# angepasst)}]
src/
|- p0a_arm_interfaces/
|  |- srv/SetNamedPose.srv  # neu
|  |- CMakeLists.txt        # angepasst: rosidl_generate_interfaces(srv)
|  `- package.xml           # angepasst: rosidl_default_runtime + member_of_group
`- p0a_arm/
   |- setup.py                     # angepasst: entry_points -> named_pose_server
   |- package.xml                  # angepasst: <depend>p0a_arm_interfaces</depend>
   `- p0a_arm/named_pose_server.py # neu
\end{lstlisting}

\subsection{Schritt 7 -- Inverse Kinematik allgemein: Denavit-Hartenberg + Jacobi}
Für zwei Glieder gäbe es eine geschlossene Lösung (Kosinussatz) -- doch die
\textbf{lässt sich nicht auf komplexe Roboter übertragen}. Reale Arme (6\,DoF und
mehr) löst man \emph{numerisch}: Man beschreibt die Kinematik mit
\textbf{Denavit-Hartenberg-Parametern} $(a_i,\alpha_i,d_i,\theta_i)$, bildet die
Vorwärtskinematik als Produkt der Gliedtransformationen und iteriert die IK über
die \textbf{Jacobi-Matrix}. Genau das tun auch \code{KDL} und \code{TRAC-IK}
hinter MoveIt\,2. Wir schreiben es so, dass derselbe Code von 2 auf $N$ Gelenke
skaliert.

\paragraph{DH-Tabelle des Arms.}
\begin{center}
\begin{tabular}{@{}ccccc@{}}
\toprule
$i$ & $a_i$ & $\alpha_i$ & $d_i$ & $\theta_i$ \\
\midrule
1 & $l_1$ & $0$ & $0$ & $\theta_1$ (Gelenk) \\
2 & $l_2$ & $0$ & $0$ & $\theta_2$ (Gelenk) \\
\bottomrule
\end{tabular}
\end{center}

\paragraph{Vorwärtskinematik.} Jede Zeile liefert eine homogene Transformation;
das Produkt ergibt die Endeffektor-Lage:
\[
A_i=\begin{bmatrix}
c\theta_i & -s\theta_i c\alpha_i & s\theta_i s\alpha_i & a_i c\theta_i\\
s\theta_i & c\theta_i c\alpha_i & -c\theta_i s\alpha_i & a_i s\theta_i\\
0 & s\alpha_i & c\alpha_i & d_i\\
0 & 0 & 0 & 1
\end{bmatrix},\qquad
T^0_n=\prod_{i=1}^{n} A_i,\qquad
\mathbf{p}(\boldsymbol\theta)=T^0_n[\,0\!:\!3,\,3\,].
\]

\paragraph{Numerische IK (Damped Least Squares).} Mit dem Fehler
$\mathbf{e}=\mathbf{p}^\ast-\mathbf{p}(\boldsymbol\theta)$ und der Jacobi-Matrix
$J=\partial\mathbf{p}/\partial\boldsymbol\theta$ iteriert man bis
$\lVert\mathbf{e}\rVert$ klein ist:
\[
\Delta\boldsymbol\theta = J^\top\!\big(JJ^\top+\lambda^2 I\big)^{-1}\mathbf{e},
\qquad \boldsymbol\theta \leftarrow \boldsymbol\theta+\Delta\boldsymbol\theta .
\]
Die Dämpfung $\lambda$ hält die Lösung auch nahe Singularitäten stabil.
\begin{lstlisting}[style=py,caption={\code{ik.py} -- DH-Vorwärtskinematik + numerische IK (skaliert auf N Gelenke)}]
import numpy as np

def dh(a, alpha, d, theta):
    ct, st, ca, sa = np.cos(theta), np.sin(theta), np.cos(alpha), np.sin(alpha)
    return np.array([[ct, -st*ca,  st*sa, a*ct],
                     [st,  ct*ca, -ct*sa, a*st],
                     [0,      sa,     ca,    d ],
                     [0,       0,      0,    1 ]])

# DH-Tabelle: (a, alpha, d) je Gelenk; theta ist die Gelenkvariable
DH = [(1.0, 0.0, 0.0), (0.8, 0.0, 0.0)]      # l1, l2

def fk(q):
    T = np.eye(4)
    for (a, al, d), th in zip(DH, q):
        T = T @ dh(a, al, d, th)
    return T[:3, 3]                            # Endeffektor-Position (x, y, z)

def jacobian(q, eps=1e-6):                     # numerisch -> voellig allgemein
    p0 = fk(q); J = np.zeros((3, len(q)))
    for i in range(len(q)):
        dq = np.array(q, float); dq[i] += eps
        J[:, i] = (fk(dq) - p0) / eps
    return J

def ik_dls(target, q0, lam=0.1, iters=200, tol=1e-4):
    q = np.array(q0, float)
    for _ in range(iters):
        e = np.asarray(target, float) - fk(q)
        if np.linalg.norm(e) < tol:
            break
        J = jacobian(q)
        q += J.T @ np.linalg.solve(J @ J.T + lam**2*np.eye(3), e)
    return q, float(np.linalg.norm(e))         # Gelenkwinkel + Restfehler
\end{lstlisting}

\begin{warnbox}[title=Numerische IK in der Praxis]
Anders als die geschlossene Lösung hängt die numerische IK vom \textbf{Startwert}
ab (sinnvoll: der aktuelle Gelenkzustand aus \code{/joint\_states}), liefert
\emph{eine} von mehreren möglichen Lösungen und muss \textbf{Gelenklimits}
nachgelagert prüfen. Die Dämpfung $\lambda$ verhindert Ausreißer nahe
Singularitäten -- beim planaren Arm etwa, dass die $z$-Richtung gar nicht
erreichbar ist (Zielebene $z=0$). Die geschlossene 2-Glieder-Lösung eignet sich
hervorragend als \textbf{Referenz}, um die numerische Implementierung
gegenzuprüfen.
\end{warnbox}
\begin{lstlisting}[style=bash,caption={\code{ros2\_ws}-Delta nach Schritt 7 (\# neu / \# angepasst)}]
src/p0a_arm/
`- p0a_arm/ik.py  # neu (DH-Vorwaertskinematik + numerische Jacobi-IK)
\end{lstlisting}

\subsection{Schritt 8 -- Action: kartesisches Ziel anfahren}
Das Anfahren eines Punktes ist ein \emph{länger laufender} Auftrag mit
Fortschritt -- also eine \textbf{Action}.
\begin{lstlisting}[style=bash,caption={\code{action/ReachPoint.action}}]
float64 x
float64 y
---
bool success
float64 final_error
---
float64 distance_to_goal
\end{lstlisting}
\paragraph{Interface erweitern.} Actions brauchen zusätzlich \code{action\_msgs}.
Die Action neben dem \code{srv} zur Generierung hinzufügen:
\begin{lstlisting}[style=bash,caption={\code{p0a\_arm\_interfaces/CMakeLists.txt} (Ergänzung)}]
find_package(action_msgs REQUIRED)

rosidl_generate_interfaces(${PROJECT_NAME}
  "srv/SetNamedPose.srv"
  "action/ReachPoint.action"
  DEPENDENCIES action_msgs
)
\end{lstlisting}
\begin{lstlisting}[style=xml,caption={\code{p0a\_arm\_interfaces/package.xml} (Ergänzung)}]
<depend>action_msgs</depend>
\end{lstlisting}

\paragraph{Action-Server.} Er löst die IK (Startwert: aktueller Gelenkzustand
aus \code{/joint\_states}), sendet die Trajektorie an den \code{arm\_controller}
und meldet die Restdistanz als Feedback, bis das Ziel erreicht ist.
\begin{lstlisting}[style=py,caption={\code{p0a\_arm/reach\_action\_server.py} (Kern)}]
import math, time
import rclpy
from rclpy.action import ActionServer
from rclpy.callback_groups import ReentrantCallbackGroup
from rclpy.executors import MultiThreadedExecutor
from rclpy.node import Node
from builtin_interfaces.msg import Duration
from sensor_msgs.msg import JointState
from trajectory_msgs.msg import JointTrajectory, JointTrajectoryPoint
from p0a_arm_interfaces.action import ReachPoint
from p0a_arm.ik import fk, ik_dls

JOINTS = ['joint1', 'joint2']

class ReachActionServer(Node):
    def __init__(self):
        super().__init__('reach_action_server')
        cb = ReentrantCallbackGroup()
        self.joint_pos = {}
        self.pub = self.create_publisher(
            JointTrajectory, '/arm_controller/joint_trajectory', 10)
        self.create_subscription(JointState, '/joint_states',
            self.on_joint_state, 10, callback_group=cb)
        ActionServer(self, ReachPoint, '/arm/reach_point',
            execute_callback=self.execute, callback_group=cb)

    def on_joint_state(self, msg):
        for n, p in zip(msg.name, msg.position):
            self.joint_pos[n] = p

    def current_q(self):
        return [self.joint_pos.get(j, 0.0) for j in JOINTS]

    def execute(self, goal_handle):
        target = [goal_handle.request.x, goal_handle.request.y, 0.0]
        q_goal, _ = ik_dls(target, self.current_q())      # IK loesen
        traj = JointTrajectory(joint_names=JOINTS)
        pt = JointTrajectoryPoint(positions=[float(v) for v in q_goal])
        pt.time_from_start = Duration(sec=2)
        traj.points = [pt]
        self.pub.publish(traj)                             # Trajektorie senden
        fb = ReachPoint.Feedback()
        deadline = time.time() + 6.0
        while time.time() < deadline:                      # Feedback-Schleife
            ee = fk(self.current_q())
            dist = math.hypot(target[0]-ee[0], target[1]-ee[1])
            fb.distance_to_goal = float(dist)
            goal_handle.publish_feedback(fb)
            if dist < 0.02:
                break
            time.sleep(0.1)
        ee = fk(self.current_q())
        res = ReachPoint.Result()
        res.final_error = float(math.hypot(target[0]-ee[0], target[1]-ee[1]))
        res.success = res.final_error < 0.05
        goal_handle.succeed()
        return res

def main():
    rclpy.init()
    ex = MultiThreadedExecutor()
    ex.add_node(ReachActionServer())
    ex.spin()
\end{lstlisting}

\begin{warnbox}[title={MultiThreadedExecutor + ReentrantCallbackGroup}]
Der \code{execute}-Callback blockiert in seiner Feedback-Schleife. Damit die
\code{/joint\_states}-Subscription \emph{währenddessen} weiterläuft, liegen beide
in einer \code{ReentrantCallbackGroup} unter einem
\code{MultiThreadedExecutor}. Mit dem Standard-Executor bliebe der Gelenkzustand
eingefroren -- das Feedback käme nie voran.
\end{warnbox}

\noindent Einstiegspunkt in \code{setup.py}:
\begin{lstlisting}[style=py,caption={\code{p0a\_arm/setup.py} -- \code{entry\_points}}]
'reach_action_server = p0a_arm.reach_action_server:main',
\end{lstlisting}

\paragraph{Ausführen (vier Terminals).} Gazebo-Stack starten, Action-Server
bereitstellen, dann Ziele senden; RViz optional.
\begin{lstlisting}[style=bash,caption={Terminal 1 -- bauen + Gazebo-Stack}]
cd ~/ros2_ws
colcon build --packages-select p0a_arm_interfaces p0a_arm
source install/setup.bash
ros2 launch p0a_arm gazebo.launch.py
\end{lstlisting}
\begin{lstlisting}[style=bash,caption={Terminal 2 -- Action-Server}]
source ~/ros2_ws/install/setup.bash
ros2 run p0a_arm reach_action_server
\end{lstlisting}
\begin{lstlisting}[style=bash,caption={Terminal 3 -- drei kartesische Ziele anfahren (mit Feedback)}]
source ~/ros2_ws/install/setup.bash
ros2 action send_goal /arm/reach_point \
  p0a_arm_interfaces/action/ReachPoint "{x: 1.2, y: 0.4}" --feedback
ros2 action send_goal /arm/reach_point \
  p0a_arm_interfaces/action/ReachPoint "{x: 0.6, y: 0.9}" --feedback
ros2 action send_goal /arm/reach_point \
  p0a_arm_interfaces/action/ReachPoint "{x: 1.5, y: -0.3}" --feedback
\end{lstlisting}
\begin{lstlisting}[style=bash,caption={Terminal 4 (optional) -- RViz eigenständig}]
source ~/ros2_ws/install/setup.bash
rviz2 -d "$(ros2 pkg prefix --share p0a_arm)/config/arm.rviz" \
  --ros-args -p use_sim_time:=true     # mit Gazebo: sonst "No transform"
\end{lstlisting}
\noindent Alle drei Ziele liegen im erreichbaren Ring ($0{,}2 \le r \le 1{,}8$
bei $l_1=1{,}0$, $l_2=0{,}8$); der Arm fährt sie nacheinander an, \code{--feedback}
zeigt die schrumpfende Restdistanz, das Result \code{success} + \code{final\_error}.
\screenshotplaceholder{p0a-reach}{RViz (oder Gazebo), während der Arm ein
kartesisches Ziel anfährt -- am besten mit eingeblendeter Endeffektor-Spur. Belegt
die selbst implementierte IK plus den Action-Mechanismus mit Feedback.}
\begin{lstlisting}[style=bash,caption={\code{ros2\_ws}-Delta nach Schritt 8 (\# neu / \# angepasst)}]
src/
|- p0a_arm_interfaces/
|  |- action/ReachPoint.action  # neu
|  |- CMakeLists.txt            # angepasst: action + DEPENDENCIES action_msgs
|  `- package.xml               # angepasst: <depend>action_msgs</depend>
`- p0a_arm/
   |- setup.py                       # angepasst: entry_points -> reach_action_server
   `- p0a_arm/reach_action_server.py # neu
\end{lstlisting}

\section{Häufige Fehler und Checkliste (Projekt 0.A)}
\begin{warnbox}[title={Stolpersteine, die fast jeden treffen}]
\begin{itemize}
\item \textbf{Aus dem falschen Verzeichnis gestartet.} Pfade wie
  \code{src/p0a\_arm/\ldots} gelten ab \code{\textasciitilde/ros2\_ws}, nicht ab
  \code{\textasciitilde/ros2\_ws/src} -- sonst entsteht \code{src/src/\ldots}.
\item \textbf{URDF per \code{-p} an der Kommandozeile.} Das mehrzeilige XML
  bricht \code{rcl}; immer über die Launch-Datei übergeben.
\item \textbf{\code{\$(find p0a\_arm)} vor dem Bauen.} Erst \code{colcon build}
  \texttt{+} \code{source}, sonst ``package not found''.
\item \textbf{\code{urdf/}/\code{launch/}/\code{config/} fehlen in
  \code{setup.py}.} Dann liegen sie nicht im \code{install/}-Baum und
  \code{ros2 launch} findet nichts -- \code{data\_files} ergänzen, neu bauen.
\item \textbf{Terminal nicht neu gesourct.} Nach \code{apt install} oder
  \code{colcon build} ein neues Terminal öffnen oder
  \code{source \textasciitilde/.bashrc}.
\item \textbf{\code{gz} / \code{joint\_state\_publisher\_gui} nicht gefunden.}
  Paket installieren (Installationskapitel) \emph{und} ROS\,2 sourcen.
\item \textbf{Custom \code{srv}/\code{action} im Python-Paket.} Nicht möglich --
  Interfaces gehören ins \code{ament\_cmake}-Paket \code{p0a\_arm\_interfaces}.
\item \textbf{Methode wie ein Node-Member benannt} (\code{handle}, \code{context}
  \ldots). Überschreibt rclpy-Interna -- Callbacks neutral benennen
  (\code{on\_request}).
\item \textbf{Link ohne \code{<inertial>} in Gazebo.} RViz verzeiht es, Gazebo
  lässt das Glied wegfallen oder wird instabil.
\item \textbf{RViz zeigt nichts / ``Frame [map] does not exist''.} \emph{Fixed
  Frame} auf \code{world} stellen und ein \code{RobotModel}-Display
  (\code{/robot\_description}, QoS \emph{Transient Local}) hinzufügen -- oder die
  fertige \code{config/arm.rviz} laden.
\item \textbf{``No transform from [link1] to [world]'', Modell bewegt sich nicht
  (mit Gazebo).} Sim-Zeit-Problem. Zweierlei nötig: (1)~\code{/clock} von Gazebo
  nach ROS brücken (\code{ros\_gz\_bridge}, siehe \code{gazebo.launch.py}) --
  fehlt sie, warnt der \code{controller\_manager} ``No clock received''; (2)~%
  \code{robot\_state\_publisher} \emph{und} RViz mit \code{use\_sim\_time:=true}
  starten. Nur beides zusammen liefert konsistente TF-Zeitstempel.
\end{itemize}
\end{warnbox}

\begin{infobox}[title=Schnelle Prüfbefehle]
\begin{lstlisting}[style=bash]
xacro src/p0a_arm/urdf/arm.urdf.xacro > /tmp/arm.urdf   # parst die URDF?
colcon build --packages-select p0a_arm                  # baut das Paket?
ros2 pkg prefix p0a_arm                                  # wird es gefunden?
ros2 node list ; ros2 topic list                         # laufen die Knoten?
rqt_graph                                                # stimmt die Architektur?
\end{lstlisting}
\end{infobox}

\section{Was du danach beherrschst}
\begin{itemize}
\item \textbf{URDF/TF:} ein Robotermodell von Hand schreiben, in RViz prüfen.
\item \textbf{Topics:} Soll-/Ist-Größen publizieren und mitlesen.
\item \textbf{Services:} kurze Befehle mit Antwort (benannte Posen).
\item \textbf{Actions:} länger laufende Aufträge mit Feedback (Ziel anfahren).
\item \textbf{IK:} das allgemeine Verfahren -- DH-Vorwärtskinematik plus
  numerische Jacobi-IK (Damped Least Squares) -- selbst implementiert; direkt
  übertragbar auf 6-DoF-Roboter.
\item \textbf{ros2\_control + Gazebo:} dieselbe Controller-Schicht wie in den
  vier Hauptprojekten (Kapitel~\ref{ch:simreal}).
\end{itemize}
